\documentclass[11pt]{article}
\usepackage[T1]{fontenc}
\usepackage[utf8]{inputenc}
\usepackage{lmodern}
\usepackage{amsmath,amssymb,amsthm}
\usepackage[margin=26mm]{geometry}
\PassOptionsToPackage{hyphens}{url}
\usepackage[hidelinks]{hyperref}
\hypersetup{pdftitle={Unbounded coefficient oscillation after unit translation},pdfauthor={Aurantium}}
\newtheorem{theorem}{Theorem}
\newtheorem{lemma}{Lemma}
\newcommand{\Z}{\mathbb Z}
\title{Unbounded coefficient oscillation\protect\\ after unit translation}
\author{Aurantium}
\date{14 October 2026}
\begin{document}
\maketitle

\begin{abstract}
We give an explicit recursive family of polynomials with positive
nonincreasing integer coefficients whose translates by one have unbounded
coefficient oscillation. Oscillation means the number of sign changes in
the list of adjacent coefficient differences after zeros are deleted.
At recursion step $r$, the degree is $8\cdot2^r-5$ and the oscillation
count is at least $2r+1$. The proof preserves an alternating subsequence
of strict differences by an integer scaling and appends two new signs
beyond the previous degree.
\end{abstract}

\section{Statement and context}

All coefficient sequences are ordered by increasing powers of $x$.
For a polynomial $B(x)=\sum_{j=0}^{D}b_jx^j$, let $\nu(B)$ be the
number of sign changes in
\[
 (b_1-b_0,b_2-b_1,\ldots,b_D-b_{D-1})
\]
after all zero entries have been deleted. An empty list or a list with
one nonzero entry has zero sign changes. No coefficients outside
$0,\ldots,D$ are included in this definition.
A sequence is \emph{unimodal} if it weakly increases to some index and
then weakly decreases; the index may be an endpoint. Thus every
nonincreasing sequence is unimodal. Oscillation is a count of changes
of direction, not a count of peaks.

\begin{theorem}\label{thm:main}
For every finite real number $R\geq0$, there is a polynomial
$A(x)=\sum_{j=0}^{d}a_jx^j$ such that
\[
 a_0\geq a_1\geq\cdots\geq a_d\geq1,
 \qquad a_j\in\Z,
 \qquad \nu(A(x+1))\geq R.
\]
More precisely, the recursive family constructed below satisfies
\begin{equation}\label{eq:bounds}
 \deg A_r=8\cdot2^r-5,
 \qquad \nu(A_r(x+1))\geq2r+1
 \quad(r\in\Z_{\geq0}).
\end{equation}
\end{theorem}

The public transcription of a 2000 report by Alvarez, Amadis and
Rosales \cite{AAR} poses the unbounded-oscillation question for unit
translates of unimodal polynomials. It already contains a positive,
nonincreasing integer-coefficient example whose translate has two
direction changes. The point here is the uniform recursive construction,
not the existence of a first nonunimodal example. The report's authors,
title and year are also recorded in Ek's bibliography \cite{Ek}.

The direction of monotonicity matters. Boros and Moll \cite{BM} prove
unimodality after unit translation for nondecreasing nonnegative
coefficients. Alvarez et al.\ \cite{AA} extend this to positive integer
shifts, and Wang and Yeh \cite{WY} to all positive real shifts.
Those input hypotheses differ from the nonincreasing class used here.

\section{Two elementary observations}

\begin{lemma}\label{lem:signs}
Let $s$ be a finite list of real numbers. If $s$ has an ordered
subsequence of $L$ nonzero entries with alternating signs, then its
nonzero entries have at least $L-1$ sign changes.
\end{lemma}

\begin{proof}
Deleting zeros preserves the selected entries and their order.
Between each consecutive pair of selected entries, the signs begin
and end differently, so there is at least one sign change in that
interval. These $L-1$ intervals contain disjoint sets of adjacent
positions in the zero-deleted list. Their sign changes can therefore
be counted separately.
\end{proof}

For $M\in\Z_{\geq0}$, put
\begin{equation}\label{eq:SM}
 S_M(x)=\sum_{k=0}^{M}x^k.
\end{equation}
The binomial theorem and Pascal's identity give, for $0\leq j\leq M$,
\begin{align}
 [x^j]S_M(x+1)
 &=\sum_{k=j}^{M}\binom{k}{j}\notag\\
 &=\sum_{k=j}^{M}
   \left(\binom{k+1}{j+1}-\binom{k}{j+1}\right)
 =\binom{M+1}{j+1}.\label{eq:hockey}
\end{align}
Here $\binom{j}{j+1}=0$. Write
$h_j=\binom{M+1}{j+1}$ for these coefficients. Each $h_j$ lies
between $0$ and the sum $2^{M+1}$ of its binomial row. Consequently
\begin{equation}\label{eq:perturbation}
 |h_{j+1}-h_j|\leq2^{M+1}
 \qquad(0\leq j<M).
\end{equation}
All sums in this argument are finite.

\section{The recursive construction}

Set $A_0=S_3$. Once $A_r$ has degree $d_r$, define
\begin{equation}\label{eq:recursion}
 M_r=2d_r+5,\qquad C_r=2^{M_r+1}+1,
 \qquad A_{r+1}=C_rA_r+S_{M_r}.
\end{equation}
We prove simultaneously that $A_r$ has positive nonincreasing integer
coefficients and that, writing $B_r=A_r(x+1)$, its strict adjacent
differences contain an ordered alternating subsequence of length
$2r+2$, beginning positive and ending negative.

\begin{proof}[Proof of Theorem~\ref{thm:main}]
For $r=0$, the coefficients of $A_0$ are $(1,1,1,1)$, and those of
$B_0$ are $(4,6,4,1)$. Its adjacent differences are
$(2,-2,-3)$; select the first two. This proves the invariant at
$r=0$, with $d_0=3$.

Suppose the invariant holds for $r$. Abbreviate $d=d_r$, $M=M_r$,
$C=C_r$, and $A'=A_{r+1}$. If $(a_0,\ldots,a_d)$ is the
coefficient list of $A_r$, the list for $A'$ is
\begin{equation}\label{eq:input-list}
 (Ca_0+1,\ldots,Ca_d+1,1,\ldots,1).
\end{equation}
It is positive, integral and nonincreasing: the old segment retains
its order, its last entry exceeds $1$, and the remaining entries are
all $1$. The leading coefficient is $1$, so $\deg A'=M$.

Write $B_r=\sum_{j=0}^{d}b_jx^j$ and
$D_j=b_{j+1}-b_j$ for $0\leq j<d$. Let
\[
 i_0<i_1<\cdots<i_{2r+1}\leq d-1
\]
be the selected indices supplied by the induction hypothesis.
Their signs are $+,-,+,-,\ldots,+,-$.
For the following coefficient calculation only, put $b_j=0$ when
$j>d$. Equation~\eqref{eq:hockey} gives the coefficients
$b'_j=Cb_j+h_j$ of $B'=A'(x+1)$ for $0\leq j\leq M$.
Thus at each old selected index,
\[
 D'_j:=b'_{j+1}-b'_j=CD_j+(h_{j+1}-h_j).
\]
The old coefficients $b_j$ are integers, because $A_r\in\Z[x]$.
Every selected $D_j$ is therefore a nonzero integer. By
\eqref{eq:perturbation},
\begin{equation}\label{eq:strict-survival}
 |CD_j|\geq C>2^{M+1}\geq|h_{j+1}-h_j|.
\end{equation}
Hence $D'_j$ has the same strict sign as $D_j$ at every selected
position. This is where the integral lower bound on the old strict
differences is used.

Two later positions provide new signs. At $j=d+1$, both old
coefficients are zero. Since $M+1=2d+6$,
\begin{align}
 D'_{d+1}
 &=\binom{2d+6}{d+3}-\binom{2d+6}{d+2}\notag\\
 &=\frac{1}{d+3}\binom{2d+6}{d+2}>0,\label{eq:new-positive}
\end{align}
where the ratio of the first binomial coefficient to the second is
$(d+4)/(d+3)>1$.
At $j=M-1$, both old coefficients again vanish, and
\begin{equation}\label{eq:new-negative}
 D'_{M-1}=\binom{M+1}{M+1}-\binom{M+1}{M}
 =1-(M+1)=-M<0.
\end{equation}
The indices occur in the required order:
\begin{equation}\label{eq:index-order}
 i_{2r+1}\leq d-1<d+1<M-1=2d+4.
\end{equation}
Appending $d+1$ and $M-1$ therefore extends the selected signs
by $+,-$, to a sequence of length $2r+4$. Both are actual
adjacent differences within the new degree $M$; no artificial
endpoint differences are counted. This completes the induction.

Lemma~\ref{lem:signs} now gives $\nu(B_r)\geq2r+1$.
The exact degree follows from the leading coefficient in
\eqref{eq:input-list}:
\[
 d_0=3,\qquad d_{r+1}=2d_r+5,
 \qquad d_r+5=8\cdot2^r.
\]
Thus \eqref{eq:bounds} holds. Finally, for any finite $R\geq0$,
choose
\begin{equation}\label{eq:R-choice}
 r=\max\left\{0,\left\lceil\frac{R-1}{2}\right\rceil\right\}.
\end{equation}
Then $2r+1\geq R$, and $A=A_r$ has every required property.
\end{proof}

\section{Scope and source provenance}

The result gives a lower bound on oscillation, not an exact count
for every $r$ or a minimum possible degree. The input monotonicity
is weak and has an endpoint mode. The statement uses the fixed
shift $1$ and does not assert a simultaneous conclusion for all
positive real shifts. The proof is symbolic and self-contained;
finite examples are not premises for its unbounded conclusion.

The report mirror \cite{AAR} displays the date 10 August 2000,
the oscillation question in Section 3 on printed page 99, and a
definition in terms of changes of direction. It does not explicitly
specify how zero differences are treated. Our zero-deletion convention
is stated above. An institutional original or page image of that
report, and its complete subsequent resolution history, have not
been authenticated here. Ek's reference \cite{Ek} corroborates the
bibliographic identity, not that history. Accordingly, no claim of
first discovery, first nonunimodal example, or worldwide priority
is made.

\paragraph{Acknowledgement.}
This note was prepared with AI assistance.

\begin{thebibliography}{9}\small
\bibitem{AAR}
J.~Alvarez, M.~Amadis and L.~Rosales,
\emph{Unimodality and Log-Concavity of Polynomials},
Summer Institute in Mathematics for Undergraduates (SIMU), technical
report, 2000. Public transcription dated 10 August 2000, Section 3,
printed p.~99; institutional original not authenticated here.
\url{https://www.academia.edu/121303677/Unimodality_and_Log_Concavity_of_Polynomials}.

\bibitem{Ek}
B.~Ek, \emph{q-Binomials and related symmetric unimodal polynomials},
author manuscript, 16 April 2018. Reference [AAR00], p.~32.
\url{https://sites.math.rutgers.edu/~bte14/Articles/Gnk/Gnk.pdf}.

\bibitem{BM}
G.~Boros and V.~H.~Moll, A criterion for unimodality,
\emph{Electron. J. Combin.} \textbf{6} (1999), R10.
\url{https://doi.org/10.37236/1442}.

\bibitem{AA}
J.~Alvarez, M.~Amadis, G.~Boros, D.~Karp, V.~H.~Moll and L.~Rosales,
An extension of a criterion for unimodality,
\emph{Electron. J. Combin.} \textbf{8} (2001), R30.
\url{https://doi.org/10.37236/1574}.

\bibitem{WY}
Y.~Wang and Y.-N.~Yeh, Proof of a conjecture on unimodality,
\emph{European J. Combin.} \textbf{26} (2005), 617--627.
Official author version: \url{https://arxiv.org/abs/0809.1586v1}.
\end{thebibliography}
\end{document}
